% SPDX-License-Identifier: MIT
% Part of lgpsf — https://github.com/NickAlger/lgpsf
%
% Mathematical companion to dev/pencil-corrections-plan.md: what the
% corrections layer computes, from what ingredients, by what formulas —
% including the one deviation from the original plan (the two-coefficient
% mode block) and the contradictions that forced it.
%
% Build: pdflatex pencil-corrections-notes.tex  (twice, for references)
\documentclass[11pt]{article}

\usepackage[margin=1.1in]{geometry}
\usepackage{amsmath, amssymb}
\usepackage{booktabs}
\usepackage[colorlinks=true, linkcolor=blue, urlcolor=blue]{hyperref}

\newcommand{\Bt}{\widetilde{B}}
\newcommand{\Hd}{H_d}
\newcommand{\Hr}{H_r}
\newcommand{\Ccorr}{C_{\mathrm{corr}}}
\newcommand{\Csurr}{C_{\mathrm{surr}}}
\newcommand{\ip}[2]{\left\langle #1, #2 \right\rangle_{\Hr}}
\newcommand{\floor}{\lambda_{\mathrm{floor}}}
\newcommand{\R}{\mathbb{R}}

\title{The pencil corrections layer:\\ operations, ingredients, and formulas}
\author{lgpsf dev notes}
\date{August 1, 2026}

\begin{document}
\maketitle

\noindent
These notes are the mathematical companion to
\texttt{dev/pencil-corrections-plan.md}. They explain, from the basics, what
the corrections layer computes, from what ingredients, and by what formulas
--- and they explain the one deviation from the original plan (the mode
block carries \emph{two} coefficient matrices rather than one) by exhibiting
the contradictions that forced it. Everything here is implemented in
\texttt{include/lgpsf/corrections/} through the flip
(slices 1--5 of the plan); deflation (slices 8--9) is included so the
operation set is complete, and is marked as planned.

\section{The problem and the ingredients}

We have fitted a sparse, symmetric operator $\Bt \in \R^{N \times N}$ that
approximates a data-misfit Hessian $\Hd$ (for lgpsf, $\Bt$ is the assembled
fit with \texttt{Symmetrize::Weighted}; validated at field scale on the Pine
Island Glacier ice-sheet Hessian). The consumer wants to solve
regularized systems
\begin{equation}
  (\Hd + a\,\Hr)\, x = b,
  \label{eq:target}
\end{equation}
where $\Hr$ is their symmetric positive definite (SPD) regularization
operator and $a > 0$ is a regularization parameter that \emph{varies}: an
L-curve sweep reuses one Hessian at many values of $a$, so every operation
below takes $a$ as an argument, and the build-time shift $a_0$ appears only
in the contracts.

The fit $\Bt$ is not directly usable in \eqref{eq:target}, for two reasons.
First, it is indefinite: each row is fitted independently, so $\Bt$ carries
a tail of small spurious negative eigenvalues even though $\Hd$ is positive
semidefinite. Second, it is inexact: the fit compresses $k \sim 100$ probe
applications of $\Hd$, and the residual error is significant in a
predictable place (directions the probes resolved but the local model could
not represent). The corrections layer repairs both --- PSD-ification for the
first, deflation for the second --- and packages the result in two
deployable forms (Section~\ref{sec:two-operators}).

\medskip
\noindent
The layer is \emph{operator-blind}: nothing in it reads a matrix entry.
The complete list of ingredients is:
\begin{enumerate}
  \item \textbf{apply $\Bt$} --- a caller-supplied symmetric operator,
    dimension plus block matvec (\texttt{SymmetricOp});
  \item \textbf{apply $\Hr$} --- cheap (mass and stiffness assemblies);
  \item \textbf{solve $\Hr$ to a tolerance} --- the consumer's
    \emph{oracle} (\texttt{HrOracle}), in practice Krylov preconditioned by
    multigrid; this is the expensive primitive, and every spectral
    operation costs $O(\#\text{modes})$ of them and nothing larger;
  \item \textbf{the probe archive} --- the pairs $(Z, Y)$ with
    $Y = \Hd Z$ that the fit was built from, plus any held-out and
    value-pass pairs; this is the \emph{only} trace of $\Hd$ we hold;
  \item \textbf{rank-sized dense algebra} --- eigendecompositions and
    Cholesky factorizations of $\rho \times \rho$ matrices,
    $\rho \lesssim$ a few hundred.
\end{enumerate}
Optionally the consumer supplies fresh applications of $\Hd$ itself, which
unlocks the value pass. What we deliberately do \emph{not} have: matrix
entries, $N$-sized factorizations, or transpose matvecs (everything
crossing the boundary is symmetric, and the layer verifies that
stochastically at construction). The restriction is what makes the same
machinery serve a sparse assembly, a block-low-rank format, or a
matrix-free shell, and what makes a distributed-memory port a change of
vector type rather than of algorithm: every $N$-sized object below is a
tall-skinny block combined through rank-sized Gram matrices.

\section{The two operators we deploy}
\label{sec:two-operators}

All corrections are low-rank and live in one object, the \emph{mode block}
(Section~\ref{sec:block}). Writing $E$ for the block's correction and $S$
for its \emph{surrogate content} --- the distinction is the subject of
Section~\ref{sec:deviation} --- the layer deploys two different operators:
\begin{align}
  P(a) &= \Bt + E + a\,\Hr,
  \label{eq:P}\\
  M(a) &= a\,\Hr + S.
  \label{eq:M}
\end{align}
$P(a)$ is the \emph{corrected operator itself}: the full accuracy of the
sparse fit plus every correction, applied matrix-free. It is what one uses
as a proposal precision or applies inside an inner Krylov iteration; its
exact solve is iterative.

$M(a)$ is the \emph{GLR deployment object}: shift plus low rank, so its
inverse is closed-form (Section~\ref{sec:woodbury}) --- one oracle solve
and $O(N\rho)$ arithmetic per application, at \emph{every} $a$, with no
refactorization. This is the preconditioner architecture validated
end-to-end at PIG scale: the consumer's Krylov iteration on
\eqref{eq:target} is preconditioned by $M(a)^{-1}$. The division of labor
is worth stating plainly: the sparse fit is the information-compression
device, and its pencil decomposition --- obtainable from sparse applies
and oracle solves alone --- is the deployment device. $M(a)$ presents the
part of the spectrum we have resolved; $P(a)$ retains everything the fit
knows, resolved or not.

\section{Why every spectral operation is a pencil operation}
\label{sec:pencil}

The typical $\Hd$ is an integral-kernel operator whose numerical rank $r$
is fixed by the physics while the discretization dimension $N$ grows under
mesh refinement ($r \sim 10^3$ at $N \sim 6.5 \times 10^3$ for the PIG
coarse mesh, the same $r$ on refinements; $r \sim 10^4$ at $N \sim 10^6$
for a continental mesh). In the Euclidean spectrum of $\Bt$ this limit is
poison: the spectrum accumulates at $0$ \emph{from both sides}, with
$N - r$ noise eigenvalues jitter-split around zero and no principled
threshold between signal and noise. ``Flip the negative eigenvalues of
$\Bt$'' is an ill-posed request: it touches $\sim (N-r)/2$ modes, more on
every refinement.

The generalized \emph{pencil}
\begin{equation}
  \Bt\, v \;=\; \lambda\, \Hr\, v
  \label{eq:pencil}
\end{equation}
is the well-posed reformulation. Its eigenvalues are dimensionless
data-to-prior ratios, mesh-independent, accumulating at $0$ with only the
physics-determined informed modes away from it; equivalently, they are the
Euclidean eigenvalues of the prior-preconditioned operator
$\Hr^{-1/2} \Bt \Hr^{-1/2}$, computed without square roots. The design
rule for the whole layer follows: \emph{never enumerate modes near an
accumulation point}. ``The rightmost $k$ pencil modes'' and ``every pencil
mode below $-\gamma a_0$'' are both requests for a physics-bounded,
mesh-independent set; their Euclidean counterparts are not. Exact PSD of
$\Bt$ alone (the case $a_0 = 0$) is out of scope for the same reason: in
the $N \to \infty$ limit it requires touching unboundedly many noise
modes. Every contract below is relative to $(\Hr, \gamma a_0)$.

\section{The mode block, and the Woodbury identity that motivates it}
\label{sec:block}
\label{sec:woodbury}

Every spectral correction the layer makes is low-rank against $\Hr$, and
they all live in one object. The \emph{mode block} holds an
$\Hr$-orthonormal basis $V \in \R^{N \times \rho}$,
\begin{equation}
  V^T \Hr V = I_\rho,
  \label{eq:orthonormal}
\end{equation}
a provenance tag per column, and coefficient matrices over it. A
correction term has the form
\begin{equation}
  E \;=\; (\Hr V)\, C \,(\Hr V)^T
  \qquad\text{with $C = C^T \in \R^{\rho\times\rho}$.}
  \label{eq:blockform}
\end{equation}
The factors of $\Hr$ look odd at first sight; they are what makes the
form close under the pencil. If $(\theta_i, u_i)$ are eigenpairs of $C$,
then with $v_i = V u_i$ one checks directly, using
\eqref{eq:orthonormal}, that
\begin{equation}
  E\, v_i \;=\; \theta_i\, \Hr\, v_i,
  \label{eq:block-eigs}
\end{equation}
i.e.\ \emph{the block's pencil eigenvalues against $\Hr$ are exactly
$\operatorname{eig}(C)$}. Every positive-definiteness certificate in the
layer is analytic because of this one fact.

The second payoff is the shifted inverse. Consider
$M(a) = a\Hr + (\Hr V) C (\Hr V)^T$ as in \eqref{eq:M}, and
eigendecompose $C = U \Theta U^T$, $\Theta =
\operatorname{diag}(\theta_1, \dots, \theta_\rho)$. The Sherman--Morrison
--Woodbury formula, with $W = \Hr V$, needs the capacitance
$C^{-1} + \tfrac1a W^T \Hr^{-1} W$; but $W^T \Hr^{-1} W = V^T \Hr V = I$
by \eqref{eq:orthonormal}, so the capacitance is \emph{diagonal} in the
eigenbasis of $C$ and the inverse closes over the oracle:
\begin{equation}
  \boxed{\;
  M(a)^{-1} x
  \;=\; \tfrac1a\, \Hr^{-1} x
  \;-\; V\, U \operatorname{diag}\!\Bigl(
        \tfrac{\theta_i}{a\,(a + \theta_i)}\Bigr)\, U^T V^T x.
  \;}
  \label{eq:woodbury}
\end{equation}
One oracle solve plus $O(N\rho)$ arithmetic per application. The formula
is easiest to believe in $\Hr$-coordinates: on the eigendirection $v_i$
the right side gives $\tfrac1a - \tfrac{\theta_i}{a(a+\theta_i)} =
\tfrac1{a+\theta_i}$, which is the correct inverse of $a + \theta_i$; on
the $\Hr$-orthogonal complement of $\operatorname{span}(V)$ it gives
$\tfrac1a \Hr^{-1}$, as required. Consequently \eqref{eq:woodbury} is
exact for \emph{every} symmetric $C$, invertible or not, and
\begin{equation}
  M(a) \succ 0
  \quad\Longleftrightarrow\quad
  a > \max\bigl(0,\, -\min_i \theta_i\bigr),
  \label{eq:glr-floor}
\end{equation}
an analytic certificate (\texttt{glr\_pd\_floor}). Changing $a$ in
\eqref{eq:woodbury} is scalar arithmetic --- this is why an entire L-curve
sweep pays for no factorization anywhere.

\subsection{Merging contributions}
\label{sec:merge}

Corrections arrive incrementally (a flip pass, then a deflation, then a
deeper cache), so the block must absorb a new contribution
$(\Hr V_{\mathrm{new}})\, C_{\mathrm{new}}\, (\Hr V_{\mathrm{new}})^T$,
$V_{\mathrm{new}} \in \R^{N \times q}$, without disturbing what is already
there. Existing columns are never modified --- provenance stays
per-column, and previously computed quantities stay bitwise valid. The
merge is classical Gram--Schmidt in the $\Hr$ inner product, done with
Gram matrices so it costs $q$ oracle \emph{applies} (no solves) plus
rank-sized dense algebra:
\begin{align}
  A &= V^T \Hr V_{\mathrm{new}} \in \R^{\rho \times q}
      &&\text{(coefficients on existing columns),} \nonumber\\
  W &= V_{\mathrm{new}} - V A
      &&\text{(residual outside the block),} \nonumber\\
  G &= W^T \Hr W = Q \Lambda Q^T
      &&\text{(residual Gram, eigendecomposed),} \nonumber\\
  V_{\mathrm{add}} &= W\, Q_+ \Lambda_+^{-1/2}
      &&\text{(new $\Hr$-orthonormal columns),}
  \label{eq:merge}
\end{align}
where $Q_+ \Lambda_+$ keeps the eigenvalues above a drop tolerance
(relative to the largest candidate's squared $\Hr$-norm), and a second
orthogonalization pass makes the arithmetic trustworthy. Candidates that
lie (nearly) inside the current span produce no new columns; their
contribution lands in the coefficients instead. With
$B = V_{\mathrm{add}}^T \Hr V_{\mathrm{new}}$ and
$M = \bigl[\begin{smallmatrix} A \\ B \end{smallmatrix}\bigr]$, the
extended coefficient matrix is
\begin{equation}
  C \;\longleftarrow\;
  \begin{bmatrix} C & 0 \\ 0 & 0 \end{bmatrix}
  + M\, C_{\mathrm{new}}\, M^T ,
  \label{eq:merge-coeffs}
\end{equation}
which is exact whenever nothing was dropped: it is just
$V_{\mathrm{new}} = V A + V_{\mathrm{add}} B$ substituted into the
contribution. Merging the same directions twice therefore adds
coefficients, not columns.

\section{The spectral machine: Lanczos in the $\Hr$ inner product}
\label{sec:lanczos}

Both ends of the pencil spectrum come from one machine. The operator
$L = \Hr^{-1}(\Bt + E)$ is self-adjoint in the inner product
$\ip{x}{y} = x^T \Hr y$:
\begin{equation}
  \ip{Lx}{y} \;=\; x^T (\Bt + E)^T \Hr^{-1} \Hr\, y
  \;=\; x^T (\Bt + E)\, y \;=\; \ip{x}{Ly},
\end{equation}
so Lanczos applies verbatim with $\ip{\cdot}{\cdot}$ in place of the dot
product. One iteration costs one operator apply, one oracle solve
(applying $L$), and one $\Hr$ apply (for the norm); storing the images
$\Hr q_j$ makes every reorthogonalization projection oracle-free. We use
full reorthogonalization against the Krylov basis \emph{and against the
block}: modes already in the block are invisible to the iteration, which
is what makes every operation incremental --- nothing is ever found
twice, and a second call continues where the first stopped. For a Ritz
pair $(\theta, y = Q s)$ from the tridiagonal $T_m$, the standard residual
identity holds in the $\Hr$-norm,
\begin{equation}
  \bigl\| L y - \theta y \bigr\|_{\Hr} \;=\; \beta_m\, |s_m|,
  \label{eq:ritz-residual}
\end{equation}
so convergence is monitored for free. Two engine details found the hard
way: on breakdown (an invariant subspace), the iteration restarts with a
fresh random vector in the $\Hr$-orthogonal complement --- itself
invariant, since $L$ is self-adjoint --- and a zero coupling in $T_m$
decouples the finished block correctly; and the engine's internal seed is
salted, because a caller who built the operator's data from the same
innocently-chosen seed can otherwise hand the engine start vectors
correlated with the operator (in the adversarial case, every start vector
is an exact eigenvector and every emptiness test lies).

\texttt{extend\_modes} runs this machine and caches converged
\emph{rightmost} pairs $(\lambda_i, v_i)$ in the block; the report carries
the first value \emph{not} cached, which is the $\lambda_{k+1}$ that
governs two-level inner conditioning. \texttt{make\_pd} runs the same
machine for the leftmost end, next.

\section{PSD-ification: the $a_0$-tied flip}
\label{sec:flip}

Let $(\lambda_i, v_i)$ be pencil eigenpairs \eqref{eq:pencil} of the
current operator, and write $w_i = \Hr v_i$. With a safety factor
$\gamma \in (0,1)$ (default $\tfrac12$), we correct exactly the modes
below the threshold $-\gamma a_0$:
\begin{equation}
  \Bt_{\mathrm{pd}}
  \;=\; \Bt \;-\; c \sum_{\lambda_i < -\gamma a_0}
        \lambda_i\, w_i w_i^T,
  \qquad
  c = \begin{cases} 2 & \text{(flip)}\\ 1 & \text{(relu)}, \end{cases}
  \label{eq:flip}
\end{equation}
implemented as block entries --- the operator handle is never modified.
Using $w_i w_i^T v_i = \Hr v_i (v_i^T \Hr v_i) = \Hr v_i$, each corrected
mode moves as
\begin{equation}
  \lambda_i \;\longmapsto\; (1 - c)\,\lambda_i
  \;=\; \begin{cases} -\lambda_i & \text{(flip)}\\
                      0 & \text{(relu)}, \end{cases}
  \label{eq:flip-move}
\end{equation}
and every other mode is untouched. The slightly-negative noise tail above
the threshold is deliberately left in place --- the deployment shift
absorbs it, and enumerating it is exactly the ill-posed operation of
Section~\ref{sec:pencil}. What makes the operation a contract rather than
a heuristic is the byproduct: the same Lanczos run resolves the leftmost
\emph{surviving} pencil value $\floor \in [-\gamma a_0, 0]$, and then,
exactly,
\begin{equation}
  \Bt_{\mathrm{pd}} + a\,\Hr \;\succ\; 0
  \quad\Longleftrightarrow\quad
  a > -\floor,
  \qquad\text{and}\qquad
  \Bt_{\mathrm{pd}} + a_0 \Hr \;\succeq\; (1-\gamma)\, a_0\, \Hr.
  \label{eq:flip-contract}
\end{equation}
Certification requires the leftmost surviving Ritz pair to
\emph{converge}: Ritz values interlace the true spectrum from above, so
an unconverged leftmost estimate proves nothing. On budget exhaustion
\texttt{make\_pd} reports uncertified with progress kept --- corrected
modes stay in the block --- and a later call continues. Flip versus relu
differed by at most one PCG iteration at every depth measured on the PIG
problem; flip is the default.

\section{The deviation: one coefficient matrix cannot serve both
operators}
\label{sec:deviation}

The original plan gave the block a single coefficient matrix $C$, used
both in the correction $E$ of \eqref{eq:P} and in the surrogate $S$ of
\eqref{eq:M}. Working Section~\ref{sec:lanczos} exposed the flaw: the two
deployment objects need \emph{different numbers on the same directions}.
Two concrete contradictions, each a one-line calculation.

\medskip
\noindent
\textbf{Contradiction 1: caching a mode must not change the operator.}
For $M(a)$ to precondition well, the block must hold the operator's
rightmost pencil modes: merge a converged eigenpair $(\lambda, v)$ of
$\Bt$ with coefficient $\lambda$, so that $S$ presents it. If that same
coefficient also enters $E$, the represented operator becomes
$\Bt + \lambda\, \Hr v v^T \Hr$, and on the mode itself
\begin{equation}
  \bigl(\Bt + \lambda \Hr v v^T \Hr\bigr) v
  \;=\; \lambda \Hr v + \lambda \Hr v
  \;=\; 2\lambda\, \Hr v :
  \label{eq:double-count}
\end{equation}
caching has \emph{doubled} the mode. The cache needs coefficient
$\lambda$ in $S$ and $0$ in $E$.

\medskip
\noindent
\textbf{Contradiction 2: a flipped mode has two different values.} The
flip \eqref{eq:flip} needs coefficient $-c\lambda$ in $E$ --- that is the
surgery that moves the operator's eigenvalue to $(1-c)\lambda$. But the
surrogate must present the \emph{corrected} operator's content on that
direction, which by \eqref{eq:flip-move} is $(1-c)\lambda$. With one
matrix, either the operator is wrong (using $(1-c)\lambda$ under-corrects
$\Bt$) or the preconditioner is wrong ($M(a)$ presents the mode at
$-c\lambda = -2\lambda$ instead of $-\lambda$; off by exactly
$\lambda$).

\medskip
\noindent
The resolution is one basis, two coefficient matrices:
\begin{equation}
  E \;=\; (\Hr V)\, \Ccorr\, (\Hr V)^T,
  \qquad
  S \;=\; (\Hr V)\, \Csurr\, (\Hr V)^T,
\end{equation}
with $\Ccorr$ the correction added to $\Bt$ and $\Csurr$ the \emph{known
spectral content of $\Bt + E$} --- what $M(a)$ deploys. Sharing the basis
is right because the two matrices overlap heavily on the same directions
(a flipped mode carries both numbers), and because the merge
\eqref{eq:merge}--\eqref{eq:merge-coeffs} extends both matrices through
the same $M$. Each merge states both increments explicitly:

\begin{center}
\begin{tabular}{lcc}
\toprule
contribution & $\Delta \Ccorr$ & $\Delta \Csurr$ \\
\midrule
cache a pencil mode $(\lambda, v)$ of $\Bt + E$ & $0$ & $\lambda$ \\
plain correction $\Delta$ (e.g.\ deflation) & $\Delta$ & $\Delta$ \\
flip/relu of a mode at $\lambda$ & $-c\lambda$ & $(1-c)\lambda$ \\
\bottomrule
\end{tabular}
\end{center}

\noindent
The consistency rule behind the table: whatever a merge adds to the
operator is also known content of $\Bt + E$ and belongs in $\Csurr$;
\emph{plus} whatever content of $\Bt$ itself the caller has learned. A
flip pass knows the mode's $\Bt$-eigenvalue $\lambda$ from the same
Lanczos run, so its surrogate increment is
$\lambda + (-c\lambda) = (1-c)\lambda$ --- the table's last row is the
first two rows composed. A deflation pass knows only its correction
$\Delta$, so on its directions $\Csurr$ holds $\Delta$ but not the
(unresolved) content of $\Bt$ there; that is an acknowledged
approximation, and a later \texttt{extend\_modes} on the corrected
operator refines the surrogate exactly where it is weakest. The
invariants are pinned by tests: a cache-only merge leaves $P(a)$
value-identical while moving the floor \eqref{eq:glr-floor}, and after a
flip the surrogate presents the corrected values $-\lambda$, not
$-2\lambda$.

With the split, the certificates separate cleanly:
$\operatorname{eig}(\Ccorr)$ feeds the positive-definiteness reasoning
for $P(a)$ (the flip contract \eqref{eq:flip-contract} and, later, the
warning-zone checks for $a < a_0$), while $\operatorname{eig}(\Csurr)$
gives the exact floor \eqref{eq:glr-floor} for $M(a)$.

\section{Deflation (planned: slices 8--9)}
\label{sec:deflation}

For completeness, the remaining operations and their formulas; the
observed numbers are from the maintainer's PIG study.

\medskip
\noindent
\emph{Free deflation.}\ The archive residuals $R = Y - \Bt Z$ are exact
samples of the error action $(\Hd - \Bt) Z$, already paid for. In the
build metric $M_0 = \Bt_{\mathrm{pd}} + a_0 \Hr$, compute $M_0^{-1} R$ by
the solve path, $M_0$-orthonormalize it to $\widetilde{Q}$ via the
Cholesky factor of the Gram matrix $R^T M_0^{-1} R$ (square-root-free),
and form the small matrix
\begin{equation}
  T \;=\; \operatorname{sym}\!\Bigl(
      (\widetilde{Q}^T R)\,
      \bigl(\widetilde{Q}^T M_0 Z\bigr)^{+}_{\mathrm{rcond}}
      \Bigr),
  \label{eq:free-deflation}
\end{equation}
whose eigenpairs estimate the whitened error
$M_0^{-1/2}(\Hd - \Bt_{\mathrm{pd}})M_0^{-1/2}$ on the residual subspace.
The truncated pseudoinverse is essential: with the raw pseudoinverse the
estimated values reached $27.7$ against a certified extreme of $3.10$ and
the corrected operator lost definiteness. Kept modes are clamped below at
the clamp floor (default $-0.95$) and merged as a plain correction. The
residual \emph{subspace} is accurate (principal cosines $0.99$ against
certified directions); the \emph{values} are the limiting factor ---
iteration counts improved $34 \to 22$ at the best setting and no setting
helped at $k \le 20$.

\medskip
\noindent
\emph{Value pass.}\ If the consumer will spend $m$ fresh applications of
$\Hd$, spend them on values, not directions: take the top-$m$ left
singular basis $Q$ of the whitened residuals (Gram-based), compute
$\Hd Q$, and form the exact Rayleigh matrix of the whitened error on that
basis. The new pairs $(Q, \Hd Q)$ enter the archive as secant
information, reusable by any later rebuild. On PIG this was decisive:
$k{=}100$ fit plus a $50$-apply value pass gave 5/12/18/24 iterations (to
$10^{-2}/10^{-4}/10^{-6}/10^{-8}$) against the true operator, versus
12/26/36/45 undeflated, with zero clamp events.

\section{Status}

Implemented and tested through slice 5: the boundary
(\texttt{symmetric\_op.hpp}, \texttt{hr\_oracle.hpp}), the two-matrix
mode block \eqref{eq:blockform}/\eqref{eq:merge-coeffs}
(\texttt{mode\_block.hpp}), the struct with the Woodbury deployment
\eqref{eq:woodbury}/\eqref{eq:glr-floor}
(\texttt{shifted\_operator.hpp}), and the spectral machine with
\texttt{extend\_modes} and \texttt{make\_pd}
\eqref{eq:flip}--\eqref{eq:flip-contract}
(\texttt{pencil\_lanczos.hpp}). Remaining: zone semantics for $a < a_0$,
the two-level solve, deflation \eqref{eq:free-deflation}, the value pass,
\texttt{rebuild\_at}, the capability-gated Cholesky backend, and the
maintainer-side validation runs. The plan of record is
\texttt{dev/pencil-corrections-plan.md}; decisions are logged in
\texttt{dev/design-notes.md}.

\end{document}
